\documentclass[journal]{IEEEtran}

% ---------------------------------------------------------------------
% Packages
% ---------------------------------------------------------------------
\usepackage{amsmath,amsfonts,amssymb}
\usepackage{algorithmic}
\usepackage{array}
\usepackage[caption=false,font=normalsize,labelfont=sf,textfont=sf]{subfig}
\usepackage{textcomp}
\usepackage{stfloats}
\usepackage{url}
\usepackage{graphicx}
\usepackage{booktabs}
\usepackage{multirow}
\usepackage{xcolor}
\usepackage[T1]{fontenc}
\usepackage{tikz}
\usetikzlibrary{trees,positioning,shapes.geometric,arrows.meta,fit,backgrounds}
\usepackage[edges]{forest}
\usepackage{cite}
\usepackage{balance}

\hyphenation{op-tical net-works semi-conduc-tor IEEE-Xplore}

% ---------------------------------------------------------------------
% Colors used by the taxonomy figure
% ---------------------------------------------------------------------
\definecolor{axisA}{RGB}{31,119,180}
\definecolor{axisB}{RGB}{44,160,44}
\definecolor{axisC}{RGB}{148,103,189}
\definecolor{rootgray}{RGB}{80,80,80}

\begin{document}

\title{A Survey on Reinforcement Learning for \\ LLM Training}

\author{Gustavo~Pantuza,
        Daniel~Neves,
        Nat\'alia~Sales,
        Caio~Campos,
        Edwaldo~Rodrigues
        
\thanks{The authors are with the Department of Computer Science,
Universidade Federal de Minas Gerais (UFMG), Belo Horizonte, Brazil
(e-mail: \{pantuza, daniel.neves, natalia.sales, caio.campos, edwaldo.rodrigues\}@dcc.ufmg.br).}
\thanks{Manuscript prepared June 2026.}}

\markboth{IEEE Communications Surveys \& Tutorials,~Vol.~XX, No.~X, 2026}%
{Campos \MakeLowercase{\textit{et al.}}: A Survey on Reinforcement Learning for LLM Training}

\maketitle

\begin{abstract}
Reinforcement learning (RL) has become the central mechanism for
transforming large language models (LLMs) from passive next-token
predictors into capable problem solvers and autonomous decision-making
agents. From the alignment of instruction-following assistants with
human preferences to the emergence of long-horizon reasoning in
models trained with verifiable rewards, RL now spans the entire
post-training stage of the modern LLM lifecycle. However, the field
has grown rapidly and unevenly, with overlapping terminology, diverse
algorithms, and heterogeneous reward designs that make systematic
comparison difficult. This survey provides a comprehensive, tutorial
review of RL for LLM training. We first formalize the paradigm shift
from single-step \emph{preference-based reinforcement fine-tuning}
(PBRFT) to temporally extended \emph{agentic RL} through a unified
Markov decision process lens, and recall the algorithmic spectrum
spanning REINFORCE, PPO, DPO, and GRPO. We then introduce a compact
yet expressive three-axis taxonomy that organizes the literature by
\emph{(i)} the reward signal that drives learning, \emph{(ii)} the
optimization algorithm that updates the policy, and \emph{(iii)} the
training scope that determines whether a single response or a full
multi-step trajectory is optimized. Using this taxonomy, we map
representative works, contrast their trade-offs, and consolidate the
landscape of open challenges, including reward hacking, credit
assignment, sample efficiency, and scalable oversight. We conclude by
outlining research directions toward reliable and general-purpose
agentic intelligence.
\end{abstract}

\begin{IEEEkeywords}
Large language models, reinforcement learning, RLHF, RLAIF,
direct preference optimization, GRPO, reward modeling, LLM agents,
alignment, taxonomy.
\end{IEEEkeywords}

% =====================================================================
% 01-introduction.tex  (IEEE Conference format — no subsections)
% =====================================================================
\section{Introduction}\label{sec:introduction}

\IEEEPARstart{R}{einforcement} Learning (RL) has become the defining
mechanism of the Large Language Model (LLM) post-training pipeline.
After pre-training by next-token prediction and optional Supervised
Fine-Tuning (SFT), it is RL that transforms a fluent yet undirected text
generator into a helpful, harmless, and increasingly \emph{agentic}
system. Reinforcement Learning from Human Feedback
(RLHF)~\cite{christiano2017deeprl,ouyang2022instructgpt} enabled the
first generation of instruction-following assistants; Direct Preference
Optimization (DPO)~\cite{rafailov2023dpo} demonstrated that the same
alignment objective could be achieved without an explicit reward model;
and Group Relative Policy Optimization (GRPO) with verifiable rewards
produced models whose reasoning capabilities are substantially shaped by
RL~\cite{shao2024deepseekmath,guo2025deepseekr1}. Most recently, this
machinery has been extended beyond single responses to optimize
multi-step \emph{trajectories} in which an LLM plans, invokes tools,
maintains memory, and acts in dynamic
environments~\cite{zhang2026landscape}.

This rapid progress, however, has outpaced systematization. Terminology
overlaps across paradigms, algorithms proliferate with variants such as
Proximal Policy Optimization (PPO), DPO, GRPO, and REINFORCE
Leave-One-Out (RLOO), and reward designs range from scalar human
preferences to dense step-level verifiers. Several surveys address parts
of this landscape: some focus on RLHF
alignment~\cite{kaufmann2024survey,casper2023open,lambert2024rlhfbook},
others catalogue broad agent
architectures~\cite{wang2024agentsurvey,xi2025agentic,masterman2024landscape},
and a recent comprehensive taxonomy organizes agentic RL by capabilities
and task domains~\cite{zhang2026landscape}. Yet none provides a
\emph{compact, tutorial-oriented} synthesis that places any given method
along a small set of clearly defined axes. Practitioners entering the
field still lack a concise map connecting preference-based fine-tuning to
the current frontier of agentic RL.

This survey fills that gap with a unified treatment of RL for LLM
training, from its inception in preference-based fine-tuning to
multi-step agentic optimization, organized around a novel three-axis
taxonomy. Our contributions are as follows:

\begin{itemize}
\item A unified Markov Decision Process (MDP) formalization of the
      paradigm shift from single-step \emph{Preference-Based
      Reinforcement Fine-Tuning} (PBRFT) to temporally extended
      \emph{agentic RL}, together with a summary of the algorithmic
      spectrum spanning REINFORCE, PPO, DPO, and GRPO.

\item A \emph{compact three-axis taxonomy} that classifies methods by
      (i)~reward signal, (ii)~optimization algorithm, and
      (iii)~training scope (single response vs.\ multi-step
      trajectory). The three axes are nearly orthogonal, so every method
      occupies a unique, well-defined cell.

\item A structured review of representative works along each axis,
      covering reward design, optimization algorithms, and agentic
      trajectories, consolidated in a master placement table.

\item A structured review of representative works along each axis,
      covering reward design, optimization algorithms, and agentic
      trajectories, consolidated in a master placement table.

\item A synthesis of open challenges, including reward hacking, credit
      assignment, sample efficiency, scalable oversight, and evaluation,
      with concrete research directions.
\end{itemize}

\noindent The remainder of this paper is organized as follows.
Section~\ref{sec:background} formalizes the RL framework for LLMs and
recalls core algorithms. Section~\ref{sec:taxonomy} introduces the
three-axis taxonomy. Sections~\ref{sec:reward}, \ref{sec:optimization},
and~\ref{sec:agentic} examine each axis in turn.
Section~\ref{sec:challenges} discusses open problems, and
Section~\ref{sec:conclusion} concludes.

% =====================================================================
% 02-background.tex
% =====================================================================
\section{Background: From LLM RL to Agentic RL}\label{sec:background}

\subsection{The LLM Post-training Pipeline}

An LLM is first \emph{pre-trained} by maximum-likelihood next-token
prediction on large corpora, yielding a general-purpose conditional
generator. Two post-training stages then shape its behavior. In
SFT the model imitates curated
$(\text{prompt},\text{response})$ demonstrations; this is stable but
bounded by the quality of the demonstrations and offers no mechanism to
exceed them. Reinforcement Fine-Tuning (RFT) instead optimizes
the model against a reward signal, allowing it to explore the response
space and, in principle, surpass its demonstrators~\cite{ouyang2022instructgpt}.

The canonical RLHF pipeline of InstructGPT chains the two: SFT, then
training a reward model on pairwise human preferences, then policy
optimization with PPO under a Kullback--Leibler (KL) penalty that keeps
the policy close to the SFT reference~\cite{ouyang2022instructgpt,ziegler2019finetuning,stiennon2020summarize}. Notably, a 1.3B-parameter
model aligned this way was preferred by humans over the 175B GPT-3,
establishing that alignment via RL can be more effective than scaling
parameters alone.

\subsection{A Unified MDP/POMDP View}

Following Zhang \textit{et~al.}~\cite{zhang2026landscape}, we cast both
classical and agentic RL for LLMs as Partially Observable Markov Decision
Processes (POMDPs), which exposes their essential difference. An MDP is a
tuple $\langle \mathcal{S}, \mathcal{A}, P, R, T, \gamma \rangle$ with
state space $\mathcal{S}$, action space $\mathcal{A}$, transition kernel
$P$, reward $R$, horizon $T$, and discount $\gamma$.

\paragraph{PBRFT.} Classical RLHF/DPO is a
\emph{degenerate, single-step} MDP. The episode starts from one prompt
state $s_0$ and terminates immediately after the model emits a complete
response $a$:
\begin{equation}
\langle \mathcal{S}_{\text{trad}}, \mathcal{A}_{\text{trad}},
P_{\text{trad}}, R_{\text{trad}}, T{=}1, \gamma{=}1 \rangle ,
\end{equation}
with $\mathcal{S}_{\text{trad}}=\{\text{prompt}\}$, a deterministic
transition to the terminal state, and a single scalar reward
$R_{\text{trad}}(s_0,a)=r(a)$. The objective is simply
$J(\theta)=\mathbb{E}_{a\sim\pi_\theta}[r(a)]$.

\paragraph{Agentic RL.} When the LLM acts over many steps, reading tool
outputs, observing an environment, and updating memory, the process becomes
a temporally extended POMDP
$\langle \mathcal{S}_{\text{agent}}, \mathcal{A}_{\text{agent}},
P_{\text{agent}}, R_{\text{agent}}, \gamma, \mathcal{O} \rangle$ with
horizon $T>1$. The agent receives observations $o_t=\mathcal{O}(s_t)$,
chooses actions $a_t$, and the state evolves stochastically
$s_{t+1}\sim P(s_{t+1}\mid s_t,a_t)$. The action space splits into
language and environment actions,
$\mathcal{A}_{\text{agent}}=\mathcal{A}_{\text{text}}\cup\mathcal{A}_{\text{action}}$, where structured actions (delimited by
special tokens) invoke tools or modify the environment. Rewards become
step-wise, mixing sparse task completion with dense sub-rewards, and the
objective is the discounted return
\begin{equation}
J(\theta)=\mathbb{E}_{\tau\sim\pi_\theta}\Big[\textstyle\sum_{t=0}^{T-1}
\gamma^t R_{\text{agent}}(s_t,a_t)\Big] .
\end{equation}

Table~\ref{tab:pbrft-vs-agentic} summarizes the contrast. The shift from
$T{=}1$ to $T{>}1$ is the single most important conceptual move in the
field: it introduces planning, long-horizon credit assignment, and
exploration as first-class concerns.

\begin{table}[t]
\centering
\caption{Formal comparison of PBRFT and Agentic RL.}
\label{tab:pbrft-vs-agentic}
\renewcommand{\arraystretch}{1.25}
\begin{tabular}{@{}p{0.10\columnwidth}p{0.36\columnwidth}p{0.40\columnwidth}@{}}
\toprule
& PBRFT & Agentic RL \\
\midrule
State & $\{s_0\}$, single prompt; ends at once & $s_t\in\mathcal{S}$, $o_t=\mathcal{O}(s_t)$, $T>1$ \\
Action & Text sequence & $\mathcal{A}_{\text{text}}\cup\mathcal{A}_{\text{action}}$ \\
Transition & Deterministic, terminal & $P(s_{t+1}\mid s_t,a_t)$ \\
Reward & Scalar $r(a)$ & Step-wise; sparse + dense \\
Objective & $\mathbb{E}[r(a)]$ & $\mathbb{E}[\sum_t \gamma^t R(s_t,a_t)]$ \\
\bottomrule
\end{tabular}
\end{table}

\subsection{The Algorithmic Spectrum}\label{sec:background-algos}

Four algorithm families recur throughout the survey; we introduce them
here and analyze their trade-offs in Section~\ref{sec:optimization}.

\paragraph{REINFORCE.} The foundational policy
gradient~\cite{williams1992reinforce,sutton2018rl} increases the
log-probability of actions weighted by their baseline-subtracted
return,
$\nabla_\theta J=\mathbb{E}[(R(s_0,a)-b(s_0))\nabla_\theta\log\pi_\theta(a\mid s_0)]$.
It is simple but suffers from high gradient variance.

\paragraph{PPO.} PPO~\cite{schulman2017ppo}
stabilizes policy-gradient updates with a clipped surrogate objective and
a learned value function (critic) for advantage estimation
$A(s_t,a_t)=R(s_t,a_t)-V(s_t)$. PPO became the workhorse of RLHF but
requires a critic of comparable size to the policy.

\paragraph{DPO.} DPO~\cite{rafailov2023dpo}
removes the explicit reward model entirely, reformulating the
KL-constrained reward-maximization as a binary classification loss over
preferred/dispreferred pairs $(y_w,y_l)$, with the policy itself acting
as an implicit reward model.

\paragraph{GRPO.} GRPO~\cite{shao2024deepseekmath}
removes the critic: it samples a \emph{group} of responses per prompt and
normalizes their rewards within the group to form advantages. Combined
with verifiable rewards, GRPO drove the emergent reasoning of
DeepSeek-R1~\cite{guo2025deepseekr1}.

% =====================================================================
% 03-taxonomy.tex  (core contribution)
% =====================================================================
\section{A Three-Axis Taxonomy}\label{sec:taxonomy}

The diversity of RL methods for LLMs can be overwhelming, but almost
every method can be characterized by answering three independent
questions: \emph{What signal does it learn from? How is the policy
updated? And over what temporal scope?} We therefore organize the field
along three nearly orthogonal axes, illustrated in
Fig.~\ref{fig:taxonomy}.

\begin{figure*}[t]
\centering
\begin{forest}
forked edges,
for tree={grow'=0,draw,rounded corners,font=\footnotesize,align=left,
          edge+={semithick},parent anchor=east,child anchor=west,
          l sep=7mm,s sep=1.2mm,inner sep=2.5pt,}
[{\textbf{RL for}\\\textbf{LLM Training}}, fill=rootgray!18, align=center
  [{\textbf{Axis A}\\Reward Signal}, fill=axisA!18, align=center
    [{Human preference (RLHF)~\cite{ouyang2022instructgpt}}]
    [{AI preference (RLAIF, Constitutional AI)~\cite{bai2022constitutional}}]
    [{Implicit / data preference (DPO)~\cite{rafailov2023dpo}}]
    [{Verifiable / rule-based (RLVR)~\cite{shao2024deepseekmath}}]
    [{Granularity: outcome (ORM) vs.\ process (PRM)~\cite{lightman2023verify}}]
  ]
  [{\textbf{Axis B}\\Optimization}, fill=axisB!18, align=center
    [{Policy gradient + critic (PPO)~\cite{schulman2017ppo}}]
    [{Critic-free / group (REINFORCE, RLOO, GRPO)~\cite{shao2024deepseekmath}}]
    [{Direct / offline (DPO, IPO, KTO)~\cite{rafailov2023dpo}}]
  ]
  [{\textbf{Axis C}\\Training Scope}, fill=axisC!18, align=center
    [{Single-step / single-turn (PBRFT)~\cite{ouyang2022instructgpt}}]
    [{Multi-step / trajectory (Agentic RL)~\cite{zhang2026landscape}}]
  ]
]
\end{forest}
\caption{The proposed three-axis taxonomy for reinforcement learning in
LLM training. The three axes, reward signal (Axis~A), optimization
algorithm (Axis~B), and training scope (Axis~C), are nearly orthogonal,
so any method occupies a single cell of the resulting space.}
\label{fig:taxonomy}
\end{figure*}

\subsection{Axis A: Reward Signal}

The reward is what the policy ultimately optimizes, and its
\emph{source} and \emph{granularity} are the most consequential design
choices. By \emph{source} we distinguish: \emph{human preference}
(RLHF), where annotators rank
responses~\cite{christiano2017deeprl,ouyang2022instructgpt};
\emph{AI preference}, Reinforcement Learning from AI Feedback
(RLAIF) or Constitutional AI, where an LLM judges responses under a
written constitution or rubric~\cite{bai2022constitutional,lee2023rlaif};
\emph{implicit or data preference}, where no reward model is trained and
the preference is encoded directly into the loss, as in
DPO~\cite{rafailov2023dpo}; and \emph{verifiable or rule-based} rewards,
Reinforcement Learning with Verifiable Rewards (RLVR), where correctness
is checked by a deterministic verifier such as a math checker or unit
test~\cite{shao2024deepseekmath,guo2025deepseekr1}.

Orthogonally, by \emph{granularity} we distinguish Outcome Reward
Models (ORMs), which score only the final answer, from Process Reward
Models (PRMs), which score each intermediate
step~\cite{lightman2023verify}. Section~\ref{sec:reward} elaborates on
this axis in detail.

\subsection{Axis B: Optimization Algorithm}

Given a reward signal, the optimization algorithm determines how the
policy is updated. We group methods into three families:
\emph{policy gradient with a critic}
(PPO~\cite{schulman2017ppo}), which estimates advantages with a learned
value function; \emph{critic-free or group-relative} methods
(REINFORCE~\cite{williams1992reinforce}, RLOO~\cite{ahmadian2024rloo},
GRPO~\cite{shao2024deepseekmath}), which avoid a value network by using
baselines or within-group normalization; and \emph{direct or offline}
preference methods (DPO~\cite{rafailov2023dpo}, Identity Preference
Optimization, IPO~\cite{azar2024ipo}, Kahneman--Tversky Optimization,
KTO~\cite{ethayarajh2024kto}), which optimize a closed-form preference
loss without on-policy sampling. The mechanisms and trade-offs of each
family are analyzed in Section~\ref{sec:optimization}.

\subsection{Axis C: Training Scope}

Finally, the training scope fixes the temporal horizon over which the
reward is optimized. \emph{Single-step} methods (PBRFT) treat the task
as one prompt mapped to one response, the degenerate $T{=}1$ MDP
described in Section~\ref{sec:background}. \emph{Multi-step} methods
(Agentic RL) optimize an entire trajectory of interleaved reasoning,
tool calls, and environment actions under a POMDP with $T{>}1$,
introducing planning and long-horizon credit
assignment~\cite{zhang2026landscape}. We dedicate
Section~\ref{sec:agentic} to this emerging paradigm.

\subsection{Placing the Literature}

The power of the taxonomy is that the three axes can be read jointly:
DeepSeek-R1, for instance, is \emph{(A)}~verifiable-reward,
\emph{(B)}~GRPO, \emph{(C)}~single-step (one long reasoning response),
whereas an RL-trained software-engineering agent is
\emph{(A)}~verifiable (unit tests), \emph{(B)}~PPO or GRPO,
\emph{(C)}~multi-step.
Table~\ref{tab:placement} places representative works in this space; the
remaining sections elaborate each row.

\begin{table*}[t]
\centering
\caption{Representative RL-for-LLM methods placed on the three axes.
RLVR = Reinforcement Learning with Verifiable Rewards;
ORM/PRM = Outcome/Process Reward Model;
PBRFT = Preference-Based Reinforcement Fine-Tuning;
CoT = Chain of Thought;
TDPO = Token-level Direct Preference Optimization.}
\label{tab:placement}
\renewcommand{\arraystretch}{1.2}
\footnotesize
\begin{tabular}{@{}llllp{0.30\textwidth}@{}}
\toprule
\textbf{Method} & \textbf{Axis A: Reward} & \textbf{Axis B: Algorithm} & \textbf{Axis C: Scope} & \textbf{Domain / Note} \\
\midrule
InstructGPT~\cite{ouyang2022instructgpt}      & Human pref.\ (ORM)        & PPO            & Single-step & First RLHF assistant at scale \\
Llama~2-Chat~\cite{touvron2023llama2}          & Human pref.\ (ORM)        & PPO + reject.\ sampling & Single-step & Open RLHF at scale \\
Constitutional AI~\cite{bai2022constitutional} & AI pref.\ (ORM)           & PPO            & Single-step & Harmlessness from AI feedback \\
RLAIF~\cite{lee2023rlaif}                       & AI pref.\ (ORM)           & PPO            & Single-step & Matches RLHF on summarization \\
DPO~\cite{rafailov2023dpo}                      & Implicit pref.            & DPO (offline)  & Single-step & No explicit reward model \\
IPO~\cite{azar2024ipo}                          & Implicit pref.            & IPO (offline)  & Single-step & Theoretical fix to DPO overfitting \\
KTO~\cite{ethayarajh2024kto}                    & Implicit (unpaired)       & KTO (offline)  & Single-step & Prospect-theoretic loss \\
TDPO~\cite{zeng2024tdpo}                         & Implicit pref.            & DPO (token)    & Single-step & Token-level control \\
RLOO~\cite{ahmadian2024rloo}                     & Human pref.               & REINFORCE-leave-one-out & Single-step & Critic-free, low variance \\
DeepSeekMath~\cite{shao2024deepseekmath}        & Verifiable (ORM)          & GRPO           & Single-step & Introduces GRPO \\
DeepSeek-R1~\cite{guo2025deepseekr1}            & Verifiable (ORM)          & GRPO           & Single-step$^\dagger$ & Emergent long-CoT reasoning \\
Let's Verify~\cite{lightman2023verify}          & Verifiable (PRM)          & Search / best-of-$n$ & Single-step & Step-level supervision \\
OmegaPRM~\cite{wang2024omegaprm}                 & Verifiable (PRM)          & Auto.\ MC supervision & Single-step & Automated process labels \\
SWE-agent~\cite{yang2024sweagent}               & Verifiable (unit tests)   & PPO / GRPO     & Multi-step  & Agent--computer interface \\
ReAct~\cite{yao2023react}                        & --- (prompting)           & RL-free        & Multi-step  & Reasoning + acting scaffold \\
Toolformer~\cite{schick2023toolformer}          & Self-supervised           & SFT            & Single-step & Learns API calls \\
Voyager~\cite{wang2023voyager}                   & Environment feedback      & RL-free        & Multi-step  & Lifelong embodied agent \\
\bottomrule
\end{tabular}\\[2pt]
{\footnotesize $^\dagger$ A single long reasoning response; the
R1 pipeline also uses multi-stage RL.}
\end{table*}

% =====================================================================
% 04-reward-signals.tex
% =====================================================================
\section{The Reward Signal}\label{sec:reward}

The reward is the only place where human intent enters RL training, so
its design dominates both the quality and the failure modes of the
resulting model. We survey the four reward sources identified in our
taxonomy and the orthogonal outcome-versus-process distinction.

\subsection{Human Preference Rewards: RLHF}

RLHF learns a reward model $r_\phi$ from pairwise human comparisons,
typically with a Bradley--Terry likelihood, and then optimizes the policy
against $r_\phi$ under a KL penalty to the reference
model~\cite{christiano2017deeprl,stiennon2020summarize,ouyang2022instructgpt}.
This recipe aligned InstructGPT and was later scaled for Llama~2-Chat,
which added rejection sampling before PPO and iterative rounds of
feedback collection~\cite{touvron2023llama2}.  Human preference remains
the gold standard for subjective qualities such as helpfulness and tone,
but it is expensive, noisy across annotators, and fragile: safety
alignment can be cheaply undone by subsequent
fine-tuning~\cite{qi2024finetuning,lermen2024lora}.  Comprehensive
treatments of RLHF and its limitations are given
in~\cite{kaufmann2024survey,casper2023open,lambert2024rlhfbook}.

\subsection{AI Preference Rewards: RLAIF}

To escape the annotation bottleneck, Constitutional AI has an LLM
critique and revise its own outputs under a written set of principles,
then trains a preference model on these AI
judgments~\cite{bai2022constitutional}.  RLAIF generalizes this idea and
was shown to match RLHF on summarization and dialogue while being far
cheaper and more reproducible~\cite{lee2023rlaif}.  The trade-off is that
AI feedback inherits the judge model's biases, including verbosity,
sycophancy, and self-preference, so RLAIF and RLHF are often combined,
using humans to calibrate and AI to scale.

\subsection{Implicit / Data Preference: DPO and Variants}

DPO observes that the RLHF objective has a closed-form optimum, allowing
the reward model to be eliminated and the policy to be trained directly
on preference pairs via a classification loss~\cite{rafailov2023dpo}.
This makes alignment dramatically simpler and more stable, at the cost of
being \emph{offline}: it learns from a fixed dataset and cannot explore.
A family of refinements followed: IPO regularizes against DPO's tendency
to overfit deterministic preferences~\cite{azar2024ipo}; KTO replaces
paired data with unpaired, prospect-theoretic
signals~\cite{ethayarajh2024kto}; Token-level Direct Preference
Optimization (TDPO) moves supervision to the token level for finer
control~\cite{zeng2024tdpo}; and calibrated DPO restores
scale-comparability with explicit rewards~\cite{liu2024caldpo}.  A key
empirical finding is that preference fine-tuning benefits from
suboptimal, on-policy data~\cite{tajwar2024preference}, foreshadowing the
online methods of Section~\ref{sec:optimization}.

\subsection{Verifiable / Rule-Based Rewards: RLVR}

The most consequential recent shift is toward rewards computed by a
deterministic verifier rather than a learned model: a math answer is
either correct or not, a unit test passes or fails.  Such
\emph{verifiable} rewards are cheap, unhackable in the usual sense, and
perfectly reproducible.  DeepSeekMath introduced GRPO precisely to
exploit them at scale~\cite{shao2024deepseekmath}, and DeepSeek-R1 showed
that pure RL with verifiable rewards, even without an initial SFT stage,
induces emergent chain-of-thought, self-verification, and ``aha''
behaviors~\cite{guo2025deepseekr1}.  Verifiable rewards are now the
dominant paradigm for reasoning, math, and code, and are a major reason
RL has become a third scaling axis alongside parameters and
data~\cite{snell2024testtime,openai2024o1,muennighoff2025s1}.

\subsection{Granularity: Outcome vs.\ Process Rewards}

Cutting across all four sources is the question of \emph{when} reward is
delivered.  Outcome reward models (ORMs) score only the final answer,
providing a sparse signal that complicates credit assignment over long
reasoning chains.  Process reward models (PRMs) instead score each
intermediate step; Lightman \textit{et~al.}\ showed PRMs outperform ORMs
for verifying mathematical reasoning~\cite{lightman2023verify}.  The
obstacle is the cost of step-level labels, addressed by automated process
supervision via Monte-Carlo rollouts~\cite{wang2024omegaprm} and by
careful study of how to build PRMs robustly~\cite{luo2025prmlessons}.
Recent work questions whether the PRM advantage transfers beyond
structured, verifiable domains, leaving process supervision for
open-ended tasks an open problem.

% =====================================================================
% 05-optimization.tex
% =====================================================================
\section{The Optimization Algorithm}\label{sec:optimization}

The second axis of our taxonomy concerns how the policy is updated once a
reward is available. The dominant methods form a spectrum from general
policy gradients to specialized preference learning; their trade-offs are
summarized in Table~\ref{tab:algos}.

\subsection{Policy Gradient with a Critic: PPO}

PPO~\cite{schulman2017ppo} optimizes a clipped surrogate that prevents
destructively large updates, using a learned value function to estimate
advantages $A(s_t,a_t)=R(s_t,a_t)-V(s_t)$. It was the engine of
InstructGPT and remains a strong baseline: a careful study found that
well-tuned PPO still outperforms DPO in complex settings that require
exploration, provided the reward model, KL penalty, and batch size are
well chosen~\cite{xu2024dpoppo}. PPO's principal cost is the critic
network, of size comparable to the policy, which roughly doubles memory
and complicates training.

\subsection{Critic-Free and Group-Relative Methods}

A second family removes the critic to reduce cost and instability.
REINFORCE with a simple baseline is the classical
instance~\cite{williams1992reinforce}, and RLOO revisits it for LLMs,
using a leave-one-out average over sampled responses as a low-variance
baseline that matches or beats PPO at lower
cost~\cite{ahmadian2024rloo}. GRPO~\cite{shao2024deepseekmath} samples a
\emph{group} of responses per prompt and normalizes their rewards within
the group to form advantages, dispensing with the value network
altogether. GRPO is especially effective with verifiable binary rewards
and underpins DeepSeek-R1~\cite{guo2025deepseekr1}; subsequent analysis
characterizes its effective loss and success-amplification
dynamics~\cite{shao2025grpodynamics}, and variants further relax the
reference-model and KL requirements for stability.

\subsection{Direct and Offline Preference Optimization}

The third family, anchored by DPO~\cite{rafailov2023dpo} and its
relatives IPO~\cite{azar2024ipo} and KTO~\cite{ethayarajh2024kto},
optimizes a closed-form preference loss directly, with no on-policy
sampling and no reward model. These methods are simple and stable but
fundamentally offline, and a growing body of evidence shows that online
methods converge to better final performance because they can explore
beyond the fixed preference
distribution~\cite{tajwar2024preference,xu2024dpoppo}. This motivates
hybrid schemes such as Nash learning from human
feedback~\cite{munos2023nash} and asynchronous off-policy
RLHF~\cite{noukhovitch2025async}, which seek the stability of offline
methods with the ceiling of online ones.

\subsection{Online Versus Offline: The Central Trade-Off}

The recurring theme is the online--offline trade-off. Offline methods
(DPO family) are cheap and stable but capped by their dataset; online
methods (PPO, GRPO) explore and reach higher performance but are
computationally expensive because every update requires fresh sampling
from the full model. Critic-free online methods (GRPO, RLOO) occupy an
attractive middle ground, which largely explains their current
popularity for reasoning-oriented training.

\begin{table}[t]
\centering
\caption{Comparison of the dominant optimization algorithms.}
\label{tab:algos}
\renewcommand{\arraystretch}{1.25}
\footnotesize
\begin{tabular}{@{}p{0.18\columnwidth}p{0.16\columnwidth}p{0.13\columnwidth}p{0.13\columnwidth}p{0.20\columnwidth}@{}}
\toprule
\textbf{Method} & \textbf{Reward model} & \textbf{Critic} & \textbf{Data} & \textbf{Best for} \\
\midrule
DPO  & Implicit & No  & Offline & General alignment \\
PPO  & Explicit & Yes & Online  & Complex / exploratory tasks \\
GRPO & Explicit (verifiable) & No & Online (groups) & Verifiable reasoning \\
RLOO & Explicit & No & Online & Low-variance, cheap \\
\bottomrule
\end{tabular}
\end{table}

% =====================================================================
% 06-agentic-rl.tex
% =====================================================================
\section{From Single Responses to Agentic Trajectories}\label{sec:agentic}

The temporal scope of optimization distinguishes single-turn PBRFT from
\emph{agentic RL}, where an LLM is optimized over an entire multi-step
trajectory~\cite{zhang2026landscape}. Although the reward and
optimization dimensions were historically studied in the single-step
regime, the frontier has moved to trajectory-level training. This section
reviews that transition, the capabilities RL is used to instill, and the
task domains and infrastructure that support it.

\subsection{Why Trajectories Change the Problem}

Extending the horizon from $T{=}1$ to $T{>}1$ reintroduces the classical
difficulties of RL that PBRFT had sidestepped. Rewards become sparse and
delayed, raising the \emph{credit-assignment} problem of deciding which
of many actions caused a success or failure. The action space expands to
include tool calls and environment actions, demanding
\emph{exploration}. The policy must also maintain \emph{state}, that is,
memory of prior observations, across the episode. Process rewards
(Section~\ref{sec:reward}) and group-relative algorithms
(Section~\ref{sec:optimization}) are partial answers to these
difficulties, but they remain the defining challenges of the agentic
regime.

\subsection{RL-Optimizable Agentic Capabilities}

Following the capability view of Zhang \textit{et~al.}~\cite{zhang2026landscape},
RL is increasingly used to convert static, prompt-engineered modules into
learned, adaptive behaviors:

\begin{itemize}
\item \textit{Reasoning.} Long Chain-of-Thought (CoT), once elicited by
prompting~\cite{wei2022cot,yao2023tot,besta2024got}, is now \emph{trained}
with verifiable rewards, as in DeepSeek-R1~\cite{guo2025deepseekr1}.

\item \textit{Planning.} Plan--execute--reflect loops can be optimized
end-to-end rather than hand-designed, though LLM planning still degrades
sharply with combinatorial complexity~\cite{valmeekam2024planning}.

\item \textit{Tool use.} From self-supervised API
learning~\cite{schick2023toolformer,qin2023toolllm,patil2023gorilla},
the field is moving toward RL that learns \emph{when} and \emph{which}
tool to call based on final-answer quality.

\item \textit{Memory.} Episodic, semantic, and procedural memory are
emerging as RL-controllable policies for what to store and
retrieve~\cite{sun2025lifelong}.

\item \textit{Self-improvement.} Agents that critique and revise their
own trajectories close the loop between generation and reward.
\end{itemize}

\subsection{Task Domains}

Agentic RL has been instantiated across several domains, almost always
with verifiable rewards. In \emph{code}, SWE-agent couples an
LLM to an agent--computer interface and is evaluated on real GitHub
issues via SWE-bench~\cite{yang2024sweagent,jimenez2024swebench,wang2024codeact};
unit-test pass/fail provides a natural reward, and performance on
SWE-bench Verified rose from roughly 12.5\% to over 50\% within a year.
In \emph{math}, formal and informal verifiers drive
training~\cite{shao2024deepseekmath,lightman2023verify}. \emph{Search}
and \emph{GUI} agents optimize multi-turn retrieval and interface
manipulation. Embodied agents such as Voyager show open-ended skill
acquisition, albeit largely without RL in the loop~\cite{wang2023voyager}.

\subsection{Multi-Agent Systems}

A further extension optimizes \emph{multiple} interacting LLM agents.
Surveys catalogue debate, hierarchical orchestration, and social
simulation as the main collaboration
patterns~\cite{talebirad2024multiagent,han2025collaboration,du2023debate,hong2023metagpt,park2023generative}.
Applying RL to jointly train communication and delegation protocols among
agents remains largely unexplored and is a promising frontier. Practical
deployments of such agents, for instance intelligent agents for
computer-network configuration~\cite{reiznautt2026minicurso} and LLM
agents for communication and network
management~\cite{boateng2024llmnsm}, illustrate the demand for
trajectory-level training in realistic, tool-rich settings.

\subsection{Environments, Frameworks, and Benchmarks}

Trajectory-level training needs infrastructure: interactive environments
that expose tools and produce rewards, scalable RL frameworks that
overlap generation and learning, and benchmarks that measure multi-step
competence rather than single answers~\cite{zhang2026landscape}. The
field has rapidly accumulated open-source environments and training
frameworks, but standardized, contamination-resistant benchmarks for
\emph{adaptive, multi-turn} behavior remain scarce, a gap we return to
in Section~\ref{sec:challenges}.

% =====================================================================
% 07-challenges.tex
% =====================================================================
\section{Open Challenges and Future Directions}\label{sec:challenges}

The field's open problems fall into three broad axes: reliable reward design, efficient optimization and long-horizon scalability. These axes give rise to the specific challenges  discussed in the subsections below.

\subsection{Reward Hacking and Robust Alignment}

Learned the reward proxies remain vulnerable to exploitation, often favoring verbosity, formatting or sycophancy over genuine quality. 
Overoptimization follows predictable scaling laws~\cite{gao2023scaling}, and reward-model ensembles mitigate but do not eliminate these failures~\cite{coste2024ensembles}. Alignment is also fragile, cheap fine-tuning can undo safety training~\cite{qi2024finetuning,lermen2024lora}.
Verifiable rewards avoid many of these issues when a suiteble verifier exists. Current directions include uncertainty-calibrated reward
models, online reward-hacking detection, and information-theoretic

\subsection{Credit Assignment and Sparse Rewards}

In multi-step settings, the central difficulty is attributing a delayed,
sparse outcome to individual actions in a long trajectory. Process reward
models~\cite{lightman2023verify,wang2024omegaprm} provide denser signal,
but generating meaningful intermediate rewards without extensive
annotation, and beyond structured mathematics, remains open.
Trajectory-aware advantage estimation and automated sub-goal verification
are active research avenues.

\subsection{Sample Efficiency and Cost}

Policy optimization remains expensive under online RL becouse each update requires sampling from the full model.
Asynchronous, off-policy RLHF~\cite{noukhovitch2025async}, difficulty-targeted data selection, and rollout reuse with importance sampling all aim to reduce the number of rollouts. Bridging offline pre-training of policies with a short online
phase offers an attractive compromise between cost and performance.

\subsection{Scalable Oversight}

As models approach or exceed human ability on a task, human preference
ceases to be a reliable reward signal, since annotators cannot judge what they
cannot verify~\cite{casper2023open}. Scalable-oversight proposals include
debate, recursive task decomposition, and weak-to-strong generalization.
The broad migration toward verifiable rewards is, in part, a response to
this limitation, but most real tasks lack a crisp verifier.

\subsection{Evaluation and Benchmarks}

Evaluation has not kept pace with model capabilities. Most benchmarks focus on single-turn interactions,
while the behaviors of interest arise in is multi-turn, adaptive settings. Benchmarks are also prone to contamination as models scale. The field needs
standardized, contamination-resistant suites that measure long-horizon, tool-using competence and the robustness of alignment under distribution
shift.

% =====================================================================
% 08-conclusion.tex
% =====================================================================
\section{Conclusion}\label{sec:conclusion}

This survey organized the rapidly evolving landscape of RL-based LLM
post-training, presenting a compact taxonomy and highlighting three dominant
trends. The first is the move from subjective preferences to verifiable reward signals that are cheaper, more reproducible, and less susceptible
to hacking. The second is the convergence on critic-free, group-relative optimizers such as GRPO and RLOO that combine online exploration with the
simplicity of offline methods. The third is the expansion of training from single responses to full trajectories, which reintroduces classical
challenges of credit assignment and stateful memory for LLMs.

The open problems we surveyed, including reward hacking, credit
assignment, sample efficiency, scalable oversight, and multi-turn
evaluation, align with these trends and suggest that progress on
verifiable rewards, efficient group-relative optimization, and
trajectory-level training will be mutually reinforcing. As these
directions mature, RL will remain central to the development of reliable, general-purpose agentic capabilities.

\section*{Acknowledgment}

The authors thank the instructors and classmates in Projeto de Agentes de IA course at UFMG for discussions that shaped this survey.


\bibliographystyle{IEEEtran}
\bibliography{references}

\end{document}
